\section*{Problem 2: Benchmarks, Baselines and Evaluation Metrics for
Heterogeneous Graph Embedding}

\noindent
This answer covers graph embedding on \emph{heterogeneous} graph structures for node
classification (NC) and link prediction (LP). The homogeneous case is treated in Problem~1;
homogeneous methods appear here only in their role as type-agnostic baselines.

% =====================================================================
\subsection*{1. Set-up: HINs, network schemas and meta-paths}

\paragraph{Heterogeneous information network (HIN).}
Following Sun et al.\ (\emph{PathSim}, PVLDB 4(11):992--1003, 2011) and the survey of Shi, Li,
Zhang, Sun \& Yu (\emph{A Survey of Heterogeneous Information Network Analysis}, IEEE TKDE
29(1):17--37, 2017), a \textbf{heterogeneous information network} is a graph
\[
  G=(V,E,\phi,\psi),\qquad \phi:V\to\mathcal{A},\qquad \psi:E\to\mathcal{R},
\]
in which every node $v\in V$ carries exactly one \emph{node type} $\phi(v)$ from the type set
$\mathcal{A}$, every edge $e\in E$ carries exactly one \emph{relation type} $\psi(e)$ from
$\mathcal{R}$, and
\[
  \boxed{\;|\mathcal{A}|+|\mathcal{R}|>2\;}
\]
This inequality \emph{is} the definition of heterogeneity: when $|\mathcal{A}|=|\mathcal{R}|=1$
the network degenerates to the homogeneous graph of Problem~1. Each relation is typed at both
ends, inducing a map $A_i\xrightarrow{R}A_j$ between two specific node types; its inverse
$R^{-1}$ is in general a \emph{different} relation ($\mathsf{writes}$ vs.\ $\mathsf{written\text{-}by}$),
which is why the released benchmarks below count $6$ or $8$ ``edge types'' where a reader might
expect $3$ or $4$.

\paragraph{Network schema.}
The \textbf{network schema} $T_G=(\mathcal{A},\mathcal{R})$ is the directed graph whose vertices
are the node types $\mathcal{A}$ and whose edges are the relation types $\mathcal{R}$. It is a
meta-template of which $G$ is an instantiation: every node and edge of $G$ maps to some type in
$T_G$, so the schema states which kinds of connection the dataset can express at all. This is the
object Section 3 draws for each benchmark. Note that the drawing carries real
information here whereas it does not in the homogeneous case: Cora's schema is a single node type
with one self-relation, while DBLP's is a four-type star.

\paragraph{Meta-path.}
A \textbf{meta-path} $\mathcal{P}$ of length $\ell$ is a path defined \emph{on the schema},
\[
  \mathcal{P}\;:\;A_1\xrightarrow{R_1}A_2\xrightarrow{R_2}\cdots\xrightarrow{R_\ell}A_{\ell+1},
\]
denoting the composite relation $R=R_1\circ R_2\circ\cdots\circ R_\ell$ from type $A_1$ to type
$A_{\ell+1}$. A node sequence $v_1v_2\cdots v_{\ell+1}$ in $G$ with $\phi(v_i)=A_i$ and
$\psi(v_iv_{i+1})=R_i$ is a \emph{path instance}. Meta-paths are written by type initials, so on
DBLP $\mathrm{APA}$ abbreviates
$\text{Author}\xrightarrow{\mathsf{writes}}\text{Paper}\xrightarrow{\mathsf{written\text{-}by}}\text{Author}$.

Meta-paths matter because each defines a \emph{different} homogeneous projection of the same HIN,
and the projections carry different semantics: on DBLP $\mathrm{APA}$ means \emph{co-authorship}
while $\mathrm{APCPA}$ means \emph{published at the same venue}, and two authors may be strongly
related under one and unrelated under the other. Nearly every method in
Section 4 is a different answer to ``how should meta-path structure be
used?'' --- fixed and hand-chosen (metapath2vec, HAN), aggregated with attention (MAGNN), learned
(GTN), or discarded in favour of per-relation parameters (R-GCN, HGT, simple-HGN).

% =====================================================================
\subsection*{2. Benchmarks}

\subsubsection*{2.1 The standardised suite: HGB}

Heterogeneous GNN papers before 2021 each preprocessed DBLP/ACM/IMDB their own way, so their
reported numbers were not comparable. \textbf{HGB} (Heterogeneous Graph Benchmark; Lv, Ding, Liu
et al., KDD 2021) exists precisely to fix this: fixed splits, fixed features, one leaderboard.
Its headline statistics are Table~\ref{p2:tab:hgb}; the per-type breakdown is
Table~\ref{p2:tab:hgbdetail}.

\begin{table}[htbp]
\centering
\small
\caption{HGB (Lv et al., KDD 2021) benchmark statistics, as published.}
\label{p2:tab:hgb}
\begin{tabular}{llrrrrlc}
\toprule
Task & Dataset & \#Nodes & \#N.\,types & \#Edges & \#E.\,types & Target & \#Classes \\
\midrule
\multirow{4}{*}{NC}
 & DBLP     & 26{,}128  & 4 & 239{,}566   & 6  & author  & 4 \\
 & IMDB     & 21{,}420  & 4 & 86{,}642    & 6  & movie   & 5$^{\dagger}$ \\
 & ACM      & 10{,}942  & 4 & 547{,}872   & 8  & paper   & 3 \\
 & Freebase & 180{,}098 & 8 & 1{,}057{,}688 & 36 & book  & 7 \\
\midrule
\multirow{3}{*}{LP}
 & Amazon   & 10{,}099  & 1 & 148{,}659   & 2  & product--product & --- \\
 & LastFM   & 20{,}612  & 3 & 141{,}521   & 3  & user--artist     & --- \\
 & PubMed   & 63{,}109  & 4 & 244{,}986   & 10 & disease--disease & --- \\
\bottomrule
\end{tabular}

\vspace{2pt}
{\footnotesize $^{\dagger}$ \textbf{multi-label} --- see Section 5.}
\end{table}

\begin{table}[htbp]
\centering
\small
\caption{Per-node-type and per-relation breakdown of the four HGB node-classification
benchmarks, recomputed from the released data files (see the note below).}
\label{p2:tab:hgbdetail}
\begin{tabular}{llrrlr}
\toprule
Dataset & Node type & Count & Feat.\ dim & Relation & \#Edges \\
\midrule
\multirow{3}{*}{DBLP}
 & author $^{\star}$ & 4{,}057  & 334   & author--paper & 19{,}645 \\
 & paper             & 14{,}328 & 4{,}231 & paper--term  & 85{,}810 \\
 & term              & 7{,}723  & 50    & paper--venue  & 14{,}328 \\
 & venue             & 20       & ---   & \emph{(+3 inverse relations)} & \\
\midrule
\multirow{3}{*}{IMDB}
 & movie $^{\star}$  & 4{,}932  & 3{,}489 & movie--director & 4{,}932 \\
 & director          & 2{,}393  & 3{,}341 & movie--actor    & 14{,}779 \\
 & actor             & 6{,}124  & 3{,}341 & movie--keyword  & 23{,}610 \\
 & keyword           & 7{,}971  & ---   & \emph{(+3 inverse relations)} & \\
\midrule
\multirow{4}{*}{ACM}
 & paper $^{\star}$  & 3{,}025  & 1{,}902 & paper--cite--paper & 5{,}343 \\
 & author            & 5{,}959  & 1{,}902 & paper--ref--paper  & 5{,}343 \\
 & subject           & 56       & 1{,}902 & paper--author      & 9{,}949 \\
 & term              & 1{,}902  & ---   & paper--subject     & 3{,}025 \\
 &                   &          &       & paper--term        & 255{,}619 \\
 &                   &          &       & \emph{(+3 inverse relations)} & \\
\midrule
\multirow{4}{*}{Freebase}
 & book $^{\star}$ 40{,}402 & \multicolumn{2}{l}{film 19{,}427} & \multicolumn{2}{l}{music 82{,}351} \\
 & sports 1{,}025 & \multicolumn{2}{l}{people 17{,}641} & \multicolumn{2}{l}{location 9{,}368} \\
 & organization 2{,}731 & \multicolumn{2}{l}{business 7{,}153} & \multicolumn{2}{l}{\emph{no features at all}} \\
 & \multicolumn{5}{l}{36 typed relations (\textsc{book-and-book} 202{,}674,
   \textsc{music-and-music} 283{,}670, \dots)} \\
\bottomrule
\end{tabular}

\vspace{2pt}
{\footnotesize $^{\star}$ labelled target type. Labelled counts: DBLP 4{,}057/4{,}057 authors;
ACM 3{,}025/3{,}025 papers; IMDB 4{,}573/4{,}932 movies (of which \textbf{2{,}684 carry more than
one label}); Freebase 7{,}954/40{,}402 books. Split is $24\%$ train / $6\%$ validation / $70\%$
test on the labelled target nodes.}
\end{table}

\paragraph{Provenance of Table~\ref{p2:tab:hgbdetail}.} The HGB paper publishes only the totals of
Table~\ref{p2:tab:hgb}; it never prints a per-type breakdown. The numbers above were therefore
recomputed directly from the released \texttt{node.dat}/\texttt{link.dat}/\texttt{label.dat} files,
and every one of the four datasets reproduces its published totals exactly (nodes, edges, number of
node types, number of edge types and number of classes all match). That agreement is what licenses
the table. Two incidental findings: HGB's own \texttt{info.dat} declares the DBLP author attribute
dimension as $0$, but the data file actually carries $334$-dimensional author features --- and $334$
is exactly the feature dimension HAN publishes for DBLP, which independently corroborates the data
over the metadata; and ACM's $1{,}902$-dimensional bag-of-words is indexed by exactly the $1{,}902$
\texttt{term} nodes.

\subsubsection*{2.2 The individual benchmark graphs}

\paragraph{DBLP (bibliographic).} A computer-science bibliography extract with types
\emph{author}, \emph{paper}, \emph{term}, \emph{venue/conference}. Authors are labelled with one of
4 research areas (database, data mining, AI, information retrieval); author features are a
bag-of-words over the keywords of the papers they wrote. It is the canonical author-classification
HIN and the standard testbed for $\mathrm{APA}$ vs.\ $\mathrm{APCPA}$ semantics.

\paragraph{ACM (bibliographic).} Papers from ACM venues with types \emph{paper}, \emph{author},
\emph{subject}, \emph{term}. Papers are labelled with 3 classes (database, wireless communication,
data mining) and carry a bag-of-words feature. Distinctively, ACM has a \emph{paper--paper}
citation relation in addition to the star structure, and HGB splits it into a citation and a
reference direction --- so the schema is not a pure star.

\paragraph{IMDB (movie).} Types \emph{movie}, \emph{director}, \emph{actor} and (in HGB)
\emph{keyword}. Movies are labelled by genre. This is the benchmark where versions diverge most
sharply: MAGNN's IMDb has 3 \emph{disjoint} genres (Action/Comedy/Drama), whereas HGB's has 5
\emph{overlapping} genres, making it multi-label.

\paragraph{Freebase.} \emph{The name denotes two unrelated graphs.} HGB's Freebase is a
book-centric knowledge graph: 8 node types, 36 relations, 180{,}098 nodes, target type
\textsc{book} with 7 classes, and \textbf{no node features whatsoever} --- so it isolates how much
a model can extract from pure typed structure. HeCo's ``Freebase'' is instead a \emph{movie}
dataset (3{,}492 movies, 33{,}401 actors, 2{,}502 directors, 4{,}459 writers; target movie, 3
classes). Quoting a Freebase number without saying which one is meaningless.

\paragraph{AMiner (bibliographic).} Two scales. metapath2vec's AMiner is web-scale: 1{,}693{,}531
authors, 3{,}194{,}405 papers, 3{,}883 venues, with 8 Google-Scholar categories as labels and
246{,}678 labelled authors. HeCo's AMiner is a small subset (6{,}564 papers, 13{,}329 authors,
35{,}890 references; target paper, 4 classes). \emph{Erratum worth knowing:} metapath2vec's body
text says the dataset ``consists of 9{,}323{,}739 computer scientists'', but the authors' own
correction note states that 9{,}323{,}739 is the paper--author \emph{link} count and the author
count is 1{,}693{,}531. The wrong figure is still widely re-quoted.

\paragraph{Yelp and Douban (business / review).} HERec's Yelp has types \emph{user} (16{,}239),
\emph{business} (14{,}284), \emph{city} (47), \emph{category} (511) and \emph{compliment} (11),
with U--B 198{,}397, U--U 158{,}590, U--Co 76{,}875, B--Ci 14{,}267 and B--Ca 40{,}009 edges
(density $0.08\%$). Douban Movie has \emph{user} (13{,}367), \emph{movie} (12{,}677),
\emph{director} (2{,}449), \emph{actor} (6{,}311) and \emph{type} (38). Both are used for
rating prediction / recommendation, reported with MAE and RMSE rather than F1.

\paragraph{LastFM and MovieLens (user--item HINs).} MAGNN's Last.fm has \emph{user} (1{,}892),
\emph{artist} (17{,}632) and \emph{tag} (1{,}088), with U--U 12{,}717, U--A 92{,}834 and A--T
23{,}253; it carries \emph{no labels and no features}, so it is a pure link-prediction benchmark
(predict user--artist), with one-hot IDs as dummy inputs. MovieLens/ML-100K is used analogously as
a user--item bipartite HIN enriched with occupation/genre types. \textbf{PubMed}, in its HGB
heterogeneous form, is a 4-type biomedical graph (63{,}109 nodes, 10 relation types) used for
disease--disease link prediction --- not to be confused with the homogeneous PubMed citation
network of Problem~1.

\paragraph{OGB: \texttt{ogbn-mag} and MAG240M.} These supply the scale the older benchmarks lack.
\texttt{ogbn-mag} is a Microsoft Academic Graph subset with \emph{paper} (736{,}389),
\emph{author} (1{,}134{,}649), \emph{institution} (8{,}740) and \emph{field of study} (59{,}965)
--- 1{,}939{,}743 nodes --- and four relations: author--affiliated\_with--institution
(1{,}043{,}998), author--writes--paper (7{,}145{,}660), paper--cites--paper (5{,}416{,}271) and
paper--has\_topic--field\_of\_study (7{,}505{,}078), totalling 21{,}111{,}007 edges. The task is
predicting each paper's \textbf{venue} among \textbf{349 classes}; only \emph{paper} nodes have
features (128-d word2vec), the rest are featureless. Crucially the split is \textbf{temporal}
--- train $\le 2017$, validation 2018, test $\ge 2019$ --- not random, so it measures
generalisation forward in time. The official metric is accuracy. \textbf{MAG240M} (OGB-LSC) is the
same design two orders of magnitude larger: 121{,}751{,}666 papers, 122{,}383{,}112 authors and
25{,}721 institutions (244{,}160{,}499 nodes) with 1{,}728{,}364{,}232 edges, 768-d RoBERTa
features, and 153 arXiv subject classes.

\subsubsection*{2.3 Version incompatibility --- the caveat that matters}

The same dataset name denotes different graphs in different papers. Table~\ref{p2:tab:versions}
makes this concrete; quoting a statistic without naming the version is an error.

\begin{table}[htbp]
\centering
\small
\caption{The same dataset name denotes different graphs in different papers. Figures that
differ from the majority version are bolded.}
\label{p2:tab:versions}
\begin{tabular}{llrrrrl}
\toprule
Dataset & Version & \multicolumn{4}{c}{Node counts} & Note \\
\midrule
\multirow{3}{*}{DBLP}
 & HAN   & \multicolumn{4}{l}{paper 14{,}328 / 14{,}327, author 4{,}057, term \textbf{8{,}789}, venue 20}
         & P--T $=88{,}420$ \\
 & MAGNN / HeCo / HGB & \multicolumn{4}{l}{paper 14{,}328, author 4{,}057, term \textbf{7{,}723}, venue 20}
         & P--T $=85{,}810$ \\
 & GTN   & \multicolumn{4}{l}{18{,}405 total (= paper + author + venue; \emph{term dropped})} & 4 edge types \\
\midrule
\multirow{3}{*}{ACM}
 & HAN   & \multicolumn{4}{l}{paper 3{,}025, author 5{,}835, subject 56} & P--A $=9{,}744$ \\
 & HGB   & \multicolumn{4}{l}{paper 3{,}025, author \textbf{5{,}959}, subject 56, term 1{,}902} & P--A $=9{,}949$ \\
 & HeCo  & \multicolumn{4}{l}{paper \textbf{4{,}019}, author \textbf{7{,}167}, subject \textbf{60}} & P--A $=13{,}407$ \\
\midrule
\multirow{3}{*}{IMDB}
 & HAN   & \multicolumn{4}{l}{movie 4{,}780, director 2{,}269, actor 5{,}841} & \\
 & MAGNN & \multicolumn{4}{l}{movie 4{,}278, director 2{,}081, actor 5{,}257} & 3 classes, single-label \\
 & HGB   & \multicolumn{4}{l}{movie 4{,}932, director 2{,}393, actor 6{,}124, keyword 7{,}971}
         & 5 classes, \textbf{multi-label} \\
\bottomrule
\end{tabular}
\end{table}

\noindent Four specific discrepancies are worth flagging.
\begin{enumerate}\itemsep2pt
\item \textbf{DBLP's term count.} HAN reports 8{,}789 terms; MAGNN, HeCo and HGB all use 7{,}723.
  HGB's published total is only consistent with the latter
  ($4{,}057+14{,}328+7{,}723+20=26{,}128$, whereas the HAN figure would give $27{,}194$). HAN's own
  table is also internally inconsistent: it lists 14{,}328 papers in the Paper--Author and
  Paper--Conference rows but 14{,}327 in the Paper--Term row.
\item \textbf{HGB's ACM does not match its stated provenance.} HGB says it uses HAN's subset while
  ``preserving all edges including paper citations and references'', but its author count
  ($5{,}959$ vs.\ $5{,}835$) and P--A edge count ($9{,}949$ vs.\ $9{,}744$) differ from HAN's
  published table. Re-adding citation edges cannot change the number of \emph{authors}, so the two
  ACM graphs genuinely differ.
\item \textbf{IMDB changes task type, not just size.} HGB's IMDB is multi-label with 5 overlapping
  genres; MAGNN's is single-label with 3 disjoint ones. An F1 on one is not comparable to an F1 on
  the other even in principle (Section 5).
\item \textbf{HGB's LastFM edge count double-counts one relation.} Its node total $20{,}612$
  exactly equals MAGNN's $1{,}892+17{,}632+1{,}088$, but its edge total is $141{,}521$ against
  MAGNN's $128{,}804$. The gap is exactly $12{,}717$, the U--U count:
  $2\times12{,}717+92{,}834+23{,}253=141{,}521$. So the symmetric user--user relation is stored in
  both directions while U--A and A--T are stored in one.
\end{enumerate}

% =====================================================================
\subsection*{3. Network schemas and meta-paths}

Each schema below is $T_G=(\mathcal{A},\mathcal{R})$ for one benchmark family: circles are node
types (letter = the initial used in meta-path notation), edges are relation types. Inverse
relations are omitted for readability.

\begin{figure}[htbp]
\centering
\begin{minipage}{0.48\textwidth}\centering
\begin{tikzpicture}[scale=0.92, every node/.style={transform shape}]
  \node[p2ent, fill=\pTwoTgt, label={[p2lab]below:Paper}]  (P) at (0,0) {P};
  \node[p2ent, fill=\pTwoAgt, label={[p2lab]below:Author}] (A) at (-3.0,0) {A};
  \node[p2ent, fill=\pTwoVen, label={[p2lab]below:Venue}]  (V) at (3.3,0) {V};
  \node[p2ent, fill=\pTwoAtt, label={[p2lab]above:Term}]   (T) at (0,1.9) {T};
  \draw[p2rel] (A) -- node[p2el]{writes} (P);
  \draw[p2rel] (P) -- node[p2el]{published in} (V);
  \draw[p2rel] (P) -- node[p2el]{contains} (T);
\end{tikzpicture}
\\[2pt]{\small\textbf{DBLP}}
\end{minipage}\hfill
\begin{minipage}{0.48\textwidth}\centering
\begin{tikzpicture}[scale=0.92, every node/.style={transform shape}]
  \node[p2ent, fill=\pTwoTgt, label={[p2lab]-35:Paper}]     (P) at (0,0) {P};
  \node[p2ent, fill=\pTwoAgt, label={[p2lab]below:Author}]  (A) at (-3.0,0) {A};
  \node[p2ent, fill=\pTwoVen, label={[p2lab]below:Subject}] (S) at (3.3,0) {S};
  \node[p2ent, fill=\pTwoAtt, label={[p2lab]above:Term}]    (T) at (0,1.9) {T};
  \draw[p2rel] (A) -- node[p2el]{writes} (P);
  \draw[p2rel] (P) -- node[p2el]{belongs to} (S);
  \draw[p2rel] (P) -- node[p2el]{contains} (T);
  \draw[p2rel] (P) to[out=-60,in=-120,looseness=9] node[p2el,below]{cites / refs} (P);
\end{tikzpicture}
\\[2pt]{\small\textbf{ACM}}
\end{minipage}

\vspace{10pt}
\begin{minipage}{0.48\textwidth}\centering
\begin{tikzpicture}[scale=0.92, every node/.style={transform shape}]
  \node[p2ent, fill=\pTwoTgt, label={[p2lab]above:Movie}]    (M) at (0,0) {M};
  \node[p2ent, fill=\pTwoAgt, label={[p2lab]below:Actor}]    (A) at (-3.1,0) {A};
  \node[p2ent, fill=\pTwoVen, label={[p2lab]below:Director}] (D) at (3.4,0) {D};
  \node[p2ent, fill=\pTwoAtt, label={[p2lab]below:Keyword}]  (K) at (0,-1.9) {K};
  \draw[p2rel] (M) -- node[p2el]{acted in} (A);
  \draw[p2rel] (M) -- node[p2el]{directed by} (D);
  \draw[p2rel] (M) -- node[p2el]{has} (K);
\end{tikzpicture}
\\[2pt]{\small\textbf{IMDB}}
\end{minipage}\hfill
\begin{minipage}{0.48\textwidth}\centering
\begin{tikzpicture}[scale=0.92, every node/.style={transform shape}]
  \node[p2ent, fill=\pTwoAgt, label={[p2lab]below:User}]   (U) at (-3.1,0) {U};
  \node[p2ent, fill=\pTwoTgt, label={[p2lab]below:Artist}] (A) at (0,0) {A};
  \node[p2ent, fill=\pTwoAtt, label={[p2lab]below:Tag}]    (T) at (3.1,0) {T};
  \draw[p2rel] (U) -- node[p2el]{listens to} (A);
  \draw[p2rel] (A) -- node[p2el]{tagged as} (T);
  \draw[p2rel] (U) to[out=120,in=60,looseness=9] node[p2el,above]{friend of} (U);
\end{tikzpicture}
\\[2pt]{\small\textbf{LastFM}}
\end{minipage}

\vspace{10pt}
\begin{minipage}{0.54\textwidth}\centering
\begin{tikzpicture}[scale=0.86, every node/.style={transform shape}]
  \node[p2ent, fill=\pTwoVen, label={[p2lab]below:Institution}]    (I) at (-6.3,0) {I};
  \node[p2ent, fill=\pTwoAgt, label={[p2lab]below:Author}]         (A) at (-3.0,0) {A};
  \node[p2ent, fill=\pTwoTgt, label={[p2lab]below:Paper}]          (P) at (0,0) {P};
  \node[p2ent, fill=\pTwoAtt, label={[p2lab]below:{Field of study}}] (F) at (3.2,0) {F};
  \draw[p2rel] (A) -- node[p2el]{affiliated with} (I);
  \draw[p2rel] (A) -- node[p2el]{writes} (P);
  \draw[p2rel] (P) -- node[p2el]{has topic} (F);
  \draw[p2rel] (P) to[out=120,in=60,looseness=9] node[p2el,above]{cites} (P);
\end{tikzpicture}
\\[2pt]{\small\textbf{\texttt{ogbn-mag} / MAG240M}}
\end{minipage}\hfill
\begin{minipage}{0.44\textwidth}\centering
\begin{tikzpicture}[scale=0.86, every node/.style={transform shape}]
  \node[p2ent, fill=\pTwoTgt, label={[p2lab]above:Business}] (B) at (0,0) {B};
  \node[p2ent, fill=\pTwoAgt, label={[p2lab]below:User}]     (U) at (-3.0,0) {U};
  \node[p2ent, fill=\pTwoAtt, label={[p2lab]below:Category}] (C) at (3.3,0) {C};
  \node[p2ent, fill=\pTwoVen, label={[p2lab]below:City}]     (L) at (0,-1.9) {L};
  \node[p2ent, fill=\pTwoAtt, label={[p2lab]above:Compliment}] (O) at (-3.0,1.9) {Co};
  \draw[p2rel] (U) -- node[p2el]{reviews} (B);
  \draw[p2rel] (B) -- node[p2el]{in category} (C);
  \draw[p2rel] (B) -- node[p2el]{located in} (L);
  \draw[p2rel] (U) -- node[p2el]{receives} (O);
  \draw[p2rel] (U) to[out=150,in=210,looseness=9] node[p2el,left]{friend} (U);
\end{tikzpicture}
\\[2pt]{\small\textbf{Yelp}}
\end{minipage}

\vspace{10pt}
\begin{minipage}{0.7\textwidth}\centering
\begin{tikzpicture}[scale=0.86, every node/.style={transform shape}]
  \node[p2ent, fill=\pTwoTgt] (B) at (0,0) {B};
  \node[p2lab] at (0,-0.88) {Book};
  \node[p2ent, fill=\pTwoAtt, label={[p2lab]45:Music}]          (M)  at (45:3.0) {M};
  \node[p2ent, fill=\pTwoAgt, label={[p2lab]90:Film}]           (F)  at (90:3.0) {F};
  \node[p2ent, fill=\pTwoVen, label={[p2lab]135:Organization}]  (O)  at (135:3.0) {O};
  \node[p2ent, fill=\pTwoVen, label={[p2lab]180:People}]        (Pe) at (180:3.0) {P};
  \node[p2ent, fill=\pTwoVen, label={[p2lab]225:Location}]      (L)  at (225:3.0) {L};
  \node[p2ent, fill=\pTwoAgt, label={[p2lab]270:Sports}]        (S)  at (270:3.0) {S};
  \node[p2ent, fill=\pTwoAtt, label={[p2lab]315:Business}]      (Bu) at (315:3.0) {U};
  \draw[p2rel] (B) -- (M);  \draw[p2rel] (B) -- (F);
  \draw[p2rel] (B) -- (O);  \draw[p2rel] (B) -- (Pe);
  \draw[p2rel] (B) -- (L);  \draw[p2rel] (B) -- (S);
  \draw[p2rel] (B) -- (Bu);
  \draw[p2rel] (B) to[out=-25,in=25,looseness=10] node[p2el,right]{and} (B);
\end{tikzpicture}
\\[2pt]{\small\textbf{Freebase (HGB)} --- \textsc{book} is adjacent to all 7 other types; 36 relations in total.}
\end{minipage}
\caption{Network schemas $T_G=(\mathcal{A},\mathcal{R})$ of the principal heterogeneous
benchmarks. Circles are node types (the letter is the initial used in meta-path notation);
labelled edges are relation types. Inverse relations are omitted throughout.}
\label{p2:fig:schemas}
\end{figure}

\begin{table}[htbp]
\centering
\small
\caption{Meta-paths commonly used on each benchmark, and the semantic relation each encodes.}
\label{p2:tab:metapaths}
\begin{tabular}{llp{7.6cm}}
\toprule
Dataset & Meta-paths & Semantics \\
\midrule
DBLP & APA & two authors co-wrote a paper (co-authorship) \\
     & APCPA / APVPA & two authors published at the same venue (community) \\
     & APTPA & two authors used the same term (topical similarity) \\
\midrule
ACM  & PAP & two papers share an author \\
     & PSP & two papers share a subject \\
     & PcPAP, PrPSP & HGB variants routed through a citation (c) or reference (r) hop \\
\midrule
IMDB & MAM & two movies share an actor \\
     & MDM & two movies share a director \\
     & AMA, DMD & two actors co-starred / two directors share a movie \\
\midrule
\texttt{ogbn-mag} & PAP & two papers share an author \\
     & PAIAP & two papers from the same institution (via author--institution) \\
\midrule
LastFM & UAU & two users listen to the same artist \\
     & UATAU & two users like artists sharing a tag \\
     & AUA, ATA & two artists share a listener / a tag \\
\midrule
Yelp & UBU & two users reviewed the same business \\
     & UBCaBU & two users reviewed businesses of the same category \\
\midrule
Freebase & BB, BFB, BPB & book--book; books sharing a film / a person \\
\bottomrule
\end{tabular}
\end{table}

\clearpage
% =====================================================================
\subsection*{4. Baselines}

\subsubsection*{4.1 Homogeneous methods, reported as type-agnostic baselines}

These are reported in essentially every heterogeneous paper, and the \emph{point} of reporting
them is that they \textbf{discard $\phi$ and $\psi$ entirely} --- they run on the graph with all
type information erased (or on one meta-path-induced projection). They are the control condition:
whatever a heterogeneous model gains over them is the measured value of modelling types. HGB's
central and uncomfortable finding is that this gap is much smaller than the literature had
claimed, and that a properly tuned GAT beats most specialised HGNNs.

\begin{itemize}\itemsep2pt
\item \textbf{DeepWalk} --- Perozzi, Al-Rfou \& Skiena, KDD 2014. Truncated uniform random walks
  treated as sentences, fed to skip-gram with hierarchical softmax.
\item \textbf{LINE} --- Tang, Qu, Wang, Zhang, Yan \& Mei, WWW 2015. Explicitly preserves
  first-order (observed edge) and second-order (shared-neighbourhood) proximity, optimised with
  negative sampling and edge-sampling for large graphs.
\item \textbf{node2vec} --- Grover \& Leskovec, KDD 2016. DeepWalk with a biased second-order walk
  whose $p,q$ parameters interpolate between BFS-like (structural equivalence) and DFS-like
  (homophily) exploration.
\item \textbf{GCN} --- Kipf \& Welling, ICLR 2017. First-order spectral convolution: a layer is a
  symmetrically normalised neighbourhood average followed by a linear map and nonlinearity.
\item \textbf{GAT} --- Veli\v{c}kovi\'{c}, Cucurull, Casanova, Romero, Li\`{o} \& Bengio, ICLR
  \textbf{2018}. Replaces the fixed normalisation with multi-head self-attention over neighbours.
  (Commonly miscited as 2017, the arXiv posting date.)
\item \textbf{GraphSAGE} --- Hamilton, Ying \& Leskovec, NIPS 2017. Inductive framework that
  samples a fixed number of neighbours and aggregates with mean/LSTM/max-pool, enabling embeddings
  for unseen nodes.
\end{itemize}

\subsubsection*{4.2 Meta-path and random-walk HIN embedding}

\begin{itemize}\itemsep3pt
\item \textbf{metapath2vec} --- Dong, Chawla \& Swami, KDD 2017. Random walks are
  \emph{meta-path-guided}: at each step the walker must move to the node type prescribed by the
  meta-path scheme (e.g.\ APVPA), so the generated context is type-balanced instead of being
  dominated by the most frequent type. The walks feed a heterogeneous skip-gram. Note that the
  base model still normalises negative sampling over \emph{all} nodes regardless of type; it is the
  \textbf{metapath2vec++} variant that normalises the softmax per node type.
\item \textbf{HIN2Vec} --- Fu, Lee \& Lei, CIKM 2017. Rather than fixing meta-paths in advance, it
  learns node \emph{and} relation (meta-path) vectors jointly by training a multi-task binary
  classifier that predicts, for a node pair, which relations hold between them.
\item \textbf{ESim} --- Shang, Qu, Liu, Kaplan, Han \& Peng, arXiv:1610.09769, 2016
  (\emph{preprint only --- never published in a peer-reviewed venue}). Embeds nodes by sampling
  positive and negative path instances of user-supplied weighted meta-paths and optimising a
  noise-contrastive objective for similarity search.
\item \textbf{HERec} --- Shi, Hu, Zhao \& Yu, IEEE TKDE 31(2):357--370, \textbf{2019}. Performs
  meta-path-guided random walks, but first \emph{filters} each walk to the target type, producing
  homogeneous sequences per meta-path; the resulting embeddings are fused by a nonlinear function
  and injected into an extended matrix-factorisation recommender.
\item \textbf{PTE} --- Tang, Qu \& Mei, KDD 2015. Decomposes a heterogeneous text network
  (word--word, word--document, word--label) into bipartite graphs and jointly embeds them with the
  LINE objective, giving semi-supervised, task-tuned text embeddings.
\end{itemize}

\subsubsection*{4.3 Heterogeneous graph neural networks}

\begin{itemize}\itemsep3pt
\item \textbf{R-GCN} --- Schlichtkrull, Kipf, Bloem, van den Berg, Titov \& Welling, ESWC
  \textbf{2018}. Relation-specific message passing: each layer aggregates neighbours separately per
  relation $r$ with its own weight matrix $W_r$, normalised by $c_{i,r}=|\mathcal{N}_i^r|$, plus a
  self-loop term. Because $|\mathcal{R}|$ weight matrices overfit, it offers two \emph{alternative}
  regularisers --- \emph{basis decomposition} ($W_r$ a linear combination of $B$ shared bases) and
  \emph{block-diagonal decomposition} ($W_r$ a direct sum of small dense blocks). For link
  prediction it is used as an encoder with a \textbf{DistMult} decoder.
\item \textbf{HAN} --- Wang, Ji, Shi, Wang, Ye, Cui \& Yu, WWW 2019. Two-level attention:
  \emph{node-level} GAT-style attention within each meta-path-induced neighbour set, then
  \emph{semantic-level} attention learning a scalar importance per meta-path to fuse them. Its
  limitation --- the one MAGNN attacks --- is that it operates on meta-path-\emph{induced} graphs,
  so intermediate nodes along the meta-path are discarded.
\item \textbf{MAGNN} --- Fu, Zhang, Meng \& King, WWW 2020. Three stages: node-content
  transformation projecting differently-dimensioned type attributes into a shared space;
  \emph{intra}-meta-path aggregation that encodes the \textbf{entire path instance including
  intermediate nodes} (mean, linear, or a relational-rotation encoder borrowed from RotatE) with
  attention over instances; and \emph{inter}-meta-path semantic attention.
\item \textbf{GTN} --- Yun, Jeong, Kim, Kang \& Kim, NeurIPS 2019. \emph{Learns} meta-paths instead
  of accepting them. A $1\times1$ convolution with softmax weights forms a soft convex combination
  $Q_i$ of the per-edge-type adjacency matrices; multiplying $A=Q_1Q_2\cdots Q_l$ yields new
  meta-path adjacencies differentiably. Including the identity among the candidates lets it learn
  paths \emph{shorter} than $l$. A GCN then runs on the learned graph.
\item \textbf{HetGNN} --- Zhang, Song, Huang, Swami \& Chawla, KDD 2019. Restart-based random walks
  sample a fixed-size, type-balanced neighbour set; content of each node is encoded with a
  Bi-LSTM over its heterogeneous attributes, neighbours are aggregated per type, and types are
  combined by attention.
\item \textbf{HGT} --- Hu, Dong, Wang \& Sun, WWW 2020. Node-type-specific projections produce
  Key/Query/Value, and attention and message matrices are indexed by the \emph{meta-relation}
  $\langle\tau(s),\phi(e),\tau(t)\rangle$, with a learnable scalar prior $\mu$ per meta-relation.
  It needs no hand-designed meta-paths. Adds \emph{relative temporal encoding} for dynamic graphs
  and \emph{HGSampling} for web-scale mini-batch training.
\item \textbf{simple-HGN / HGB} --- Lv, Ding, Liu et al., KDD 2021. A deliberately minimal
  baseline: GAT plus (i) a learnable \textbf{edge-type embedding} concatenated into the attention
  logit, (ii) residual connections on \emph{both} node representations and attention scores, and
  (iii) $L_2$ normalisation of the output embeddings. The paper's larger contribution is the HGB
  benchmark itself and the finding that this simple model matches or beats most specialised HGNNs.
\item \textbf{ie-HGCN} --- Yang, Guan, Li, Zhao, Cui \& Wang, IEEE TKDE 35(2):1637--1650, 2023.
  Hierarchical object-level then type-level aggregation which, as a by-product, \emph{ranks} the
  useful meta-paths, making the model interpretable while avoiding an exhaustive meta-path search.
\item \textbf{RSHN} --- Zhu, Zhou, Pan, Zhu \& Wang, IEEE ICDM 2019. Builds a coarsened
  \emph{edge-centric} ``relation structure graph'' to embed relation types themselves, then
  propagates with those relation embeddings rather than treating relations as bare labels.
\item \textbf{HPN} --- Ji, Wang, Shi, Wang \& Yu, IEEE TKDE 35(1):521--532, 2023. Analyses
  semantic confusion (the HIN analogue of over-smoothing) and replaces the propagation step with a
  personalised-PageRank-like update that retains a fraction of the node's own representation,
  enabling deeper meta-path propagation.
\end{itemize}

\subsubsection*{4.4 Self-supervised / contrastive heterogeneous methods}

\begin{itemize}\itemsep3pt
\item \textbf{DMGI} --- Park, Kim, Han \& Yu, AAAI 2020. Extends Deep Graph Infomax to
  \emph{attributed multiplex} networks: one DGI mutual-information objective per relation type,
  with a consensus regularisation that pulls the per-relation embeddings towards a shared
  representation.
\item \textbf{HDGI} --- Ren, Liu, Huang, Dai, Bo \& Zhang, arXiv:1911.08538, 2019 (AAAI-20 DLGMA
  workshop, under a different title; \emph{no archival venue}). Applies mutual-information
  maximisation between local node embeddings and a global summary, using meta-path-based
  convolution plus semantic attention as the encoder.
\item \textbf{HeCo} --- Wang, Liu, Han \& Shi, KDD 2021. Contrasts two \emph{views} of the same
  node: a \emph{network-schema} view (attention over directly-connected neighbours of each other
  type) and a \emph{meta-path} view (meta-path-specific GCNs with semantic attention). Positives
  are not merely the node's counterpart view --- HeCo selects multiple positives by counting how
  many meta-path instances connect two nodes and taking the top ones.
\end{itemize}

\subsubsection*{4.5 Knowledge-graph embeddings as relation-aware LP baselines}

A HIN with many relation types is close to a knowledge graph, so heterogeneous link-prediction
papers (and HGB's LP track) routinely report KG embedding models, which score a typed triple
$(h,r,t)$ directly.

\begin{itemize}\itemsep2pt
\item \textbf{TransE} --- Bordes, Usunier, Garcia-Duran, Weston \& Yakhnenko, NIPS 2013. Models a
  relation as a translation, $\mathbf{h}+\mathbf{r}\approx\mathbf{t}$, scoring
  $-\|\mathbf{h}+\mathbf{r}-\mathbf{t}\|$. Cannot represent 1-to-N or symmetric relations.
\item \textbf{DistMult} --- Yang, Yih, He, Gao \& Deng, ICLR 2015. Bilinear score
  $\langle\mathbf{h},\mathrm{diag}(\mathbf{r}),\mathbf{t}\rangle$; efficient but necessarily
  \emph{symmetric}.
\item \textbf{ComplEx} --- Trouillon, Welbl, Riedel, Gaussier \& Bouchard, ICML 2016. Moves
  DistMult into $\mathbb{C}$ and takes the real part of a Hermitian product, so asymmetric
  relations become representable at the same cost.
\item \textbf{ConvE} --- Dettmers, Minervini, Stenetorp \& Riedel, AAAI 2018. Reshapes $h$ and $r$
  into 2-D and applies a convolution, giving a parameter-efficient, highly expressive scorer.
\item \textbf{RotatE} --- Sun, Deng, Nie \& Tang, ICLR 2019. Treats a relation as an element-wise
  \emph{rotation} in complex space, $\mathbf{t}=\mathbf{h}\circ\mathbf{r}$ with
  $|r_i|=1$, capturing symmetry, antisymmetry, inversion and composition simultaneously.
\end{itemize}

% =====================================================================
\subsection*{5. Evaluation metrics}

Two conventions govern the whole literature. First, a HIN benchmark normally labels
\textbf{exactly one node type} --- the author in DBLP, the paper in ACM, the movie in IMDB --- so
``accuracy on DBLP'' always means accuracy on that one target type, the other types being context.
Second, for \emph{unsupervised} methods the protocol is \textbf{embed $\to$ freeze $\to$ fit a
linear classifier}: the embedding is learned without labels, then held fixed while a logistic
regression is trained on a fraction of target nodes, the fraction varied over $20/40/60/80\%$ and
mean $\pm$ std reported over repeated splits. HAN, MAGNN and HeCo all report in this form.

\subsubsection*{5.1 Node classification}

With class set $\mathcal{C}$ and per-class counts $\mathrm{TP}_c,\mathrm{FP}_c,\mathrm{FN}_c$,
\textbf{Micro-F1} pools the confusion matrix \emph{before} scoring,
\[
  \mathrm{Micro\text{-}P}=\frac{\sum_{c}\mathrm{TP}_c}{\sum_{c}(\mathrm{TP}_c+\mathrm{FP}_c)},
  \qquad
  \mathrm{Micro\text{-}R}=\frac{\sum_{c}\mathrm{TP}_c}{\sum_{c}(\mathrm{TP}_c+\mathrm{FN}_c)},
\]
\[
  \mathrm{Micro\text{-}F1}
   =\frac{2\,\mathrm{Micro\text{-}P}\cdot\mathrm{Micro\text{-}R}}
          {\mathrm{Micro\text{-}P}+\mathrm{Micro\text{-}R}},
\]
whereas \textbf{Macro-F1} scores each class and averages \emph{unweighted}:
\[
  \mathrm{F1}_c=\frac{2\,\mathrm{TP}_c}{2\,\mathrm{TP}_c+\mathrm{FP}_c+\mathrm{FN}_c},
  \qquad
  \mathrm{Macro\text{-}F1}=\frac{1}{|\mathcal{C}|}\sum_{c\in\mathcal{C}}\mathrm{F1}_c .
\]
In the \emph{single-label} multi-class setting every misclassification creates one false positive
and one false negative, so $\mathrm{Micro\text{-}P}=\mathrm{Micro\text{-}R}$ and Micro-F1 equals
plain \textbf{accuracy}. The pair is always reported together because they answer different
questions: Micro-F1 is dominated by frequent classes, Macro-F1 weights a rare class as heavily as a
common one, and a large gap between them signals a model that has learned the class prior rather
than the graph.

\paragraph{The multi-label caveat.} This equivalence \textbf{fails} for HGB's IMDB and Freebase
tracks, where a node may carry several labels at once. Prediction is per-label thresholding (sigmoid
+ binary cross-entropy) rather than an $\arg\max$, and the pooled confusion matrix is taken over the
flattened node~$\times$~label indicator matrix; there $\mathrm{Micro\text{-}P}\neq
\mathrm{Micro\text{-}R}$ in general and Micro-F1 is \emph{not} accuracy. In HGB's IMDB, 2{,}684 of
the 4{,}573 labelled movies carry more than one genre, so this is not a corner case --- and it is
why a Micro-F1 on HGB's IMDB cannot be compared with one on MAGNN's single-label IMDB.

\paragraph{Other classification metrics.} \textbf{Accuracy} is the official metric on OGB
(\texttt{ogbn-mag}, MAG240M). \textbf{ROC-AUC} is used for binary or strongly imbalanced targets
(and by HeCo alongside F1). For the \emph{clustering} probe that HAN, MAGNN, DMGI and HeCo all run
--- $k$-means on the frozen embeddings with $k=|\mathcal{C}|$ --- the scores are
\[
  \mathrm{NMI}(Y,\hat{C})=\frac{2\,I(Y;\hat{C})}{H(Y)+H(\hat{C})}\in[0,1],
  \qquad
  \mathrm{ARI}=\frac{\mathrm{RI}-\mathbb{E}[\mathrm{RI}]}
                    {\max(\mathrm{RI})-\mathbb{E}[\mathrm{RI}]},
\]
with $I$ mutual information, $H$ entropy and $\mathrm{RI}$ the Rand index. Both are $0$ in
expectation for a chance partition and $1$ for a perfect one (ARI may go negative). These test
whether the embedding geometry is class-structured with no classifier fitted at all.

\subsubsection*{5.2 Link prediction}

The essential difference from the homogeneous case is that in a HIN link prediction is stated
\textbf{per relation type}: not ``is there an edge'' but ``is there an edge \emph{of type $R$}''
--- predict Author--Paper, or Paper--Venue, or User--Artist. Two consequences follow.
(i) \emph{Negatives must be sampled type-consistently}: a negative for an Author--Paper edge must
join an author to a paper, never an author to a venue, since a type-violating negative is trivially
rejected by any type-aware model and inflates its apparent margin over type-agnostic baselines.
HGB samples negatives from 2-hop neighbours at a 1:1 ratio for exactly this reason.
(ii) Scores are reported \emph{per relation} and then averaged, so a model can be strong on a dense
relation and useless on a sparse one without that showing in a single number.

Let $\mathrm{rank}_i$ be the position of the $i$-th true edge among its candidates, over query set
$Q$:
\[
  \mathrm{MRR}=\frac{1}{|Q|}\sum_{i\in Q}\frac{1}{\mathrm{rank}_i},
  \qquad
  \mathrm{Hits}@K=\frac{1}{|Q|}\sum_{i\in Q}\mathbf{1}\!\left[\mathrm{rank}_i\le K\right],
\]
i.e.\ MRR is the mean reciprocal rank of the first true item and Hits@$K$ the fraction of positive
edges ranked inside the top $K$ against their sampled negatives ($K\in\{1,3,10\}$ conventionally,
inherited from the KG literature). \textbf{ROC-AUC} equals the probability that a random positive
outranks a random negative. \textbf{Average precision} (equivalently AUPRC) is
\[
  \mathrm{AP}=\sum_{n}\left(R_n-R_{n-1}\right)P_n ,
\]
with $P_n,R_n$ precision and recall at the $n$-th threshold. AP is preferred to ROC-AUC when
positives are rare --- always true for link prediction on a sparse graph --- because ROC-AUC is
optimistic there, the enormous true-negative count sitting in its denominator.
\textbf{Precision@$K$} / \textbf{Recall@$K$}, and MAE/RMSE for rating prediction, appear in the
recommendation-flavoured HIN papers (HERec on Yelp/Douban). HGB's LP track standardises on
\textbf{ROC-AUC and MRR}.

\subsubsection*{5.3 Summary: dataset $\to$ target $\to$ task $\to$ metric}

\begin{table}[htbp]
\centering
\small
\caption{Dataset $\to$ target node type $\to$ task $\to$ metric normally reported.}
\label{p2:tab:summary}
\begin{tabular}{llllll}
\toprule
Dataset & Domain & Target & Task & Labels & Metrics normally reported \\
\midrule
DBLP      & bibliographic & author  & NC     & 4, single & Macro-F1, Micro-F1 (+NMI/ARI) \\
ACM       & bibliographic & paper   & NC     & 3, single & Macro-F1, Micro-F1 (+NMI/ARI) \\
IMDB (HGB)& movie         & movie   & NC     & 5, \textbf{multi} & Macro-F1, Micro-F1 \\
IMDB (MAGNN)& movie       & movie   & NC     & 3, single & Macro-F1, Micro-F1 \\
Freebase (HGB) & KG       & book    & NC     & 7, single & Macro-F1, Micro-F1 \\
AMiner    & bibliographic & author/paper & NC & 8 / 4    & Macro-F1, Micro-F1, NMI \\
\texttt{ogbn-mag} & bibliographic & paper & NC & 349, single & \textbf{Accuracy} (official) \\
MAG240M   & bibliographic & paper   & NC     & 153, single & \textbf{Accuracy} (official) \\
\midrule
LastFM    & music         & user--artist & LP & ---      & ROC-AUC, MRR \\
Amazon    & e-commerce    & product--product & LP & ---  & ROC-AUC, MRR \\
PubMed (HGB) & biomedical & disease--disease & LP & ---  & ROC-AUC, MRR \\
Yelp / Douban & review    & user--item & LP / rating & --- & Precision@K, Recall@K; MAE, RMSE \\
OAG (HGT) & bibliographic & paper--field / venue & LP & --- & NDCG, MRR \\
\bottomrule
\end{tabular}
\end{table}

\paragraph{Closing remark.} The single most important methodological point is the one HGB was
created to make: because DBLP, ACM and IMDB circulate in mutually incompatible preprocessed
versions (Table~\ref{p2:tab:versions}), and because IMDB even changes from single-label to
multi-label between versions, numbers copied across papers are frequently not comparable. Any
credible evaluation must name the dataset \emph{version}, the split, and whether the protocol is
end-to-end supervised or embed-freeze-classify.
